\documentclass[11pt]{article}
\newcommand{\DocShort}{Dossier de exposiciones}
\usepackage{fontspec}
\setmainfont{TeX Gyre Pagella}
\setsansfont{TeX Gyre Heros}
\usepackage[spanish,es-noshorthands]{babel}
\decimalpoint
\usepackage{geometry}
\geometry{letterpaper, top=2.4cm, bottom=2.4cm, left=2.7cm, right=2.7cm}
\usepackage[expansion=false]{microtype}
\usepackage{parskip}
\setlength{\parskip}{6pt}
\usepackage{amsmath,amssymb}
\usepackage{booktabs}
\usepackage{array}
\usepackage{longtable}
\usepackage{caption}
\usepackage[dvipsnames,table]{xcolor}
\usepackage{enumitem}
\setlist[itemize]{topsep=2pt, itemsep=2pt, parsep=0pt}
\setlist[enumerate]{topsep=2pt, itemsep=2pt, parsep=0pt}
\usepackage{fancyhdr}
\usepackage[hidelinks,pdfencoding=auto]{hyperref}
\usepackage{titlesec}
\usepackage{tcolorbox}
\tcbuselibrary{breakable}
\usepackage{tabularx}
\usepackage{makecell}
\usepackage{ragged2e}
\usepackage{csquotes}
\usepackage[backend=biber,style=apa,sorting=nyt,sortcites=true,language=spanish]{biblatex}
\DeclareLanguageMapping{spanish}{spanish-apa}
\addbibresource{referencias_presentaciones.bib}

\definecolor{darkblue}{HTML}{174A63}
\definecolor{medblue}{HTML}{2E7896}
\definecolor{lightblue}{HTML}{DCEEF4}
\definecolor{teal}{HTML}{1B7F79}
\definecolor{lightteal}{HTML}{D7EFEC}
\definecolor{lightgrey}{HTML}{F5F7F8}
\definecolor{midgrey}{HTML}{5B6770}
\definecolor{warnyellow}{HTML}{FFF4CC}

\titleformat{\section}
{\large\bfseries\color{darkblue}}
{\thesection}{0.75em}{}[{\color{medblue}\titlerule[0.8pt]}]
\titleformat{\subsection}
{\normalsize\bfseries\color{darkblue}}
{\thesubsection}{0.75em}{}
\titleformat{\subsubsection}
{\normalsize\bfseries\color{medblue}}
{\thesubsubsection}{0.75em}{}

\pagestyle{fancy}
\fancyhf{}
\fancyhead[L]{\small\color{midgrey}ICV-342-T Hidrología}
\fancyhead[R]{\small\color{midgrey}\DocShort}
\fancyfoot[C]{\small\color{midgrey}\thepage}
\setlength{\headheight}{14pt}
\renewcommand{\headrulewidth}{0.4pt}
\renewcommand{\footrulewidth}{0pt}

\rowcolors{2}{lightgrey}{white}
\renewcommand{\arraystretch}{1.18}
\newcolumntype{Y}{>{\RaggedRight\arraybackslash}X}

\newtcolorbox{infobox}{
  breakable,
  colback=lightblue,
  colframe=medblue,
  boxrule=0.8pt,
  arc=3pt,
  left=8pt,
  right=8pt,
  top=7pt,
  bottom=7pt
}

\newtcolorbox{warnbox}{
  breakable,
  colback=warnyellow,
  colframe=orange!70!black,
  boxrule=0.8pt,
  arc=3pt,
  left=8pt,
  right=8pt,
  top=7pt,
  bottom=7pt
}

\newcommand{\doccover}[2]{
  \begin{center}
  {\small\color{midgrey}Pontificia Universidad Católica Madre y Maestra\\Escuela de Ingeniería Civil y Ambiental}

  \vspace{1.0cm}
  {\color{medblue}\rule{\linewidth}{2pt}}
  \vspace{0.3cm}

  {\Large\bfseries\color{darkblue}ICV-342-T Hidrología}\\[0.35em]
  {\LARGE\bfseries\color{darkblue}#1}

  \vspace{0.3cm}
  {\color{medblue}\rule{\linewidth}{2pt}}
  \vspace{0.9cm}
  \end{center}

  \begin{center}
  \begin{tabular}{ll}
  \toprule
  \textbf{Asignatura} & ICV-342-T Hidrología \\
  \textbf{Profesor} & Ricardo Rafael Hernández Moreira, PhD \\
  \textbf{Periodo académico} & Mayo--agosto de 2026 \\
  \textbf{Versión} & #2 \\
  \bottomrule
  \end{tabular}
  \end{center}

  \vspace{0.4cm}
}

\newcommand{\unitbox}[3]{
  \begin{tcolorbox}[
    breakable,
    colback=lightteal,
    colframe=teal,
    boxrule=0.8pt,
    arc=3pt,
    title=\textbf{#1}
  ]
  \textbf{Enfoque de la exposición.} #2

  \vspace{0.5\baselineskip}
  \textbf{Apoyo bibliográfico sugerido.} #3
  \end{tcolorbox}
}

\AtBeginDocument{
  \hypersetup{
    pdftitle={\DocShort},
    pdfauthor={Ricardo Rafael Hernández Moreira},
    pdfsubject={ICV-342-T Hidrología}
  }
}


\begin{document}
\doccover{Dossier para coordinar exposiciones}{26 de mayo de 2026}

\begin{infobox}
Este documento organiza las exposiciones temáticas de la asignatura en las unidades II a IX. Su propósito es fijar reglas de operación, criterios de evaluación y un piso académico común para que cada equipo prepare una exposición técnicamente útil, rigurosa y comparable.
\end{infobox}

\section{Alcance y criterio general}
\begin{itemize}
\item Las exposiciones cubren las unidades II a IX del programa oficial de ICV-342-T Hidrología.
\item Cada exposición debe responder al programa de la asignatura y apoyarse, como mínimo, en \textcite{aparicio1989} y en al menos otra fuente de la carpeta de bibliografía del curso.
\item El nivel esperado es \emph{balanceado}: conceptos clave, procedimiento básico cuando aplique, interpretación ingenieril y al menos un ejemplo o ejercicio breve bien explicado.
\item La exposición no debe limitarse a leer diapositivas. El equipo debe demostrar criterio técnico, capacidad de síntesis y control del tiempo.
\end{itemize}

\section{Parámetros generales de la sesión}
\begin{center}
\begin{tabularx}{\linewidth}{p{0.25\linewidth} Y}
\toprule
\textbf{Parámetro} & \textbf{Regla de trabajo} \\
\midrule
Horario de clase & Miércoles de 17:00 a 19:00. \\
Hora límite de arranque & Las presentaciones deberán comenzar a más tardar a las 17:15. Si el equipo y el profesor están listos antes, pueden iniciar antes de esa hora. \\
Cobertura por sesión & Una unidad temática por sesión, salvo ajuste excepcional indicado por el profesor. \\
Duración nominal del equipo & Entre 35 y 40 minutos de exposición efectiva. \\
Preguntas y discusión & Después de la exposición habrá espacio para preguntas de la audiencia y cierre técnico del profesor. \\
Entrega obligatoria del equipo & Diapositivas empleadas en clase y bibliografía efectivamente utilizada. \\
Énfasis del material visual & Figuras, tablas, esquemas y ejemplos legibles. No se espera exceso de texto ni ornamentación innecesaria. \\
\bottomrule
\end{tabularx}
\end{center}

\begin{warnbox}
Una exposición que exceda el tiempo asignado perjudica la discusión técnica y el cierre docente. La versión presentada debe estar ensayada y cronometrada con anterioridad.
\end{warnbox}

\begin{warnbox}
También se penalizarán las exposiciones excesivamente cortas. Como regla operativa:
\begin{itemize}
\item Si la exposición efectiva dura entre 30 y 34 minutos, se descontará 10\% de la calificación obtenida en la rúbrica del equipo expositor.
\item Si la exposición efectiva dura entre 25 y 29 minutos, se descontará 20\% de la calificación obtenida en la rúbrica del equipo expositor.
\item Si la exposición efectiva dura menos de 25 minutos, se descontará 35\% de la calificación obtenida en la rúbrica del equipo expositor.
\end{itemize}
El objetivo de esta regla es evitar exposiciones subdesarrolladas o cerradas prematuramente. Terminar antes por pocos segundos no activa penalización; la referencia es la duración efectiva total de la exposición académica del equipo.
\end{warnbox}

\section{Expectativas académicas del equipo expositor}
\begin{itemize}
\item Presentar las definiciones fundamentales con precisión, sin quedarse en un listado de términos.
\item Explicar por qué el tema importa en hidrología aplicada y en diseño u operación de obras civiles.
\item Mostrar los procedimientos, criterios o relaciones más relevantes del tema, aunque no se desarrollen todas las derivaciones.
\item Incluir por lo menos un ejemplo, interpretación de datos, gráfico, esquema o problema breve que aterrice el contenido.
\item Cerrar con ideas operativas: qué variables se miden, qué supuestos son delicados, qué errores son frecuentes y cómo se conecta el tema con otras unidades.
\item Citar con claridad las fuentes utilizadas en el material visual y en la bibliografía final.
\item Evitar imágenes decorativas que no se comenten de forma directa durante la exposición.
\item No usar imágenes generadas por IA.
\item Citar la fuente de cada imagen o figura utilizada y preferir material de libre acceso o con permiso de uso claro.
\end{itemize}

\section{Rúbrica del equipo expositor}
La evidencia Exposición oral vale 5\% de la asignatura según el programa. Para operar estas exposiciones, se utilizará la siguiente distribución interna:

\begin{center}
\begin{tabular}{p{0.44\linewidth} p{0.18\linewidth} p{0.22\linewidth}}
\toprule
\textbf{Componente} & \textbf{Peso en el curso} & \textbf{Observación} \\
\midrule
Desempeño del equipo expositor & 4\% & Se califica con la rúbrica de esta sección. \\
Participación de audiencia y revisión breve por sesión & 1\% & Se acumula con el formulario individual de la audiencia. \\
\bottomrule
\end{tabular}
\end{center}

\small
{\setlength{\tabcolsep}{4pt}
\begin{longtable}{@{}p{0.22\linewidth} p{0.17\linewidth} p{0.17\linewidth} p{0.17\linewidth} p{0.17\linewidth}@{}}
\toprule
\textbf{Criterio} & \textbf{Excelente} & \textbf{Satisfactorio} & \textbf{En desarrollo} & \textbf{Insuficiente} \\
\midrule
\endfirsthead
\toprule
\textbf{Criterio} & \textbf{Excelente} & \textbf{Satisfactorio} & \textbf{En desarrollo} & \textbf{Insuficiente} \\
\midrule
\endhead
Cobertura del tema (20\%) & Atiende el núcleo del programa, distingue ideas centrales y no omite piezas necesarias para comprender la unidad. & Cubre la mayor parte del tema, con omisiones menores o poco impacto global. & La cobertura es parcial o desbalanceada. & Omite partes esenciales o desordena el tema de forma seria. \\
Dominio técnico (25\%) & Explica criterios, relaciones y supuestos con seguridad; responde preguntas con fundamento técnico. & Muestra dominio razonable con vacíos menores. & Se observan dudas importantes o respuestas poco precisas. & El equipo no demuestra control suficiente del contenido. \\
Claridad y estructura (15\%) & La secuencia de ideas es nítida, bien jerarquizada y fácil de seguir. & La exposición es entendible, aunque con transiciones o énfasis mejorables. & La organización complica el seguimiento del tema. & La exposición es confusa o impropia para aprendizaje del grupo. \\
Aplicación e interpretación (15\%) & Integra ejemplo, datos, gráfico o situación de ingeniería y extrae conclusiones útiles. & Incluye alguna aplicación pertinente, aunque limitada. & La aplicación aparece débil o poco conectada con el contenido. & No aterriza el contenido en un uso técnico reconocible. \\
Material visual (10\%) & Las diapositivas son legibles, limpias y técnicas; cada imagen o figura usada tiene función clara, se comenta en clase y está citada. & El material cumple, aunque con algunos problemas de legibilidad, síntesis o trazabilidad visual. & El apoyo visual dificulta parcialmente la comprensión o usa recursos gráficos poco justificados. & El material es desordenado, ilegible, improductivo o usa imágenes decorativas o no citadas. \\
Manejo del tiempo y trabajo de equipo (15\%) & El tiempo se controla bien, hay participación equilibrada y transiciones razonables. & El manejo del tiempo o la distribución del trabajo presenta debilidades menores. & Hay desequilibrio notable o problemas claros de tiempo. & La coordinación del equipo afecta seriamente la calidad de la sesión. \\
\bottomrule
\end{longtable}
}
\normalsize

\section{Evaluación de audiencia y regla de asistencia}
\begin{itemize}
\item La audiencia no es pasiva. Cada estudiante debe asistir, seguir la exposición y entregar un formulario individual breve al final de la sesión.
\item La parte de audiencia equivale a 1\% del curso, distribuida entre todas las sesiones de exposición.
\item El formulario debe llevar nombre y matrícula para que el profesor pueda registrar asistencia y calificación. Si se comparten observaciones con el equipo expositor, se hará sin identificar al estudiante que las escribió.
\item Si un estudiante falta a una sesión de exposiciones o no entrega el formulario correspondiente, obtiene cero en esa sesión para la parte de audiencia.
\item La calidad de la revisión importa: no basta con marcar presencia. Se espera al menos una observación técnica, una pregunta o una respuesta breve a la pregunta técnica planteada.
\end{itemize}

\begin{center}
\begin{tabular}{p{0.50\linewidth} p{0.34\linewidth}}
\toprule
\textbf{Criterio de audiencia} & \textbf{Se espera} \\
\midrule
Asistencia y entrega (40\%) & Presencia efectiva y entrega del formulario de la sesión. \\
Atención verificable (35\%) & Comentario, observación o pregunta conectada con el contenido expuesto. \\
Respuesta técnica breve (25\%) & Respuesta clara a una pregunta técnica corta sobre la exposición. \\
\bottomrule
\end{tabular}
\end{center}

\section{Calendario de exposiciones y exámenes}
\begin{infobox}
El calendario se organiza entre el miércoles 03 de junio de 2026 y el miércoles 29 de julio de 2026. La propuesta reserva tres días completos para pruebas escritas, dos jornadas con dos exposiciones y un día exclusivo para la Unidad VI.
\end{infobox}

\begin{center}
\begin{tabularx}{\linewidth}{p{0.16\linewidth} p{0.18\linewidth} Y}
\toprule
\textbf{Fecha} & \textbf{Formato} & \textbf{Actividad} \\
\midrule
03 jun 2026 & Exposición & Unidad II. Meteorología e hidrología. \\
10 jun 2026 & Exposición & Unidad III. Cuencas hidrográficas. \\
17 jun 2026 & Examen 100 & Examen 1 de 100 minutos sobre las unidades I, II y III. \\
24 jun 2026 & Doble exposición & Unidad IV. Precipitación, seguida de Unidad V. Escurrimiento. \\
01 jul 2026 & Exposición & Unidad VI. Predicción de eventos de diseño. Día exclusivo para esta exposición. \\
08 jul 2026 & Examen 100 & Examen 2 de 100 minutos sobre las unidades IV, V y VI. \\
15 jul 2026 & Doble exposición & Unidad VII. Evaporación y transpiración, seguida de Unidad VIII. Infiltración. \\
22 jul 2026 & Exposición & Unidad IX. Agua subterránea. \\
29 jul 2026 & Examen 100 & Examen 3 de 100 minutos sobre las unidades VII, VIII y IX. \\
\bottomrule
\end{tabularx}
\end{center}

\begin{warnbox}
En las jornadas con doble exposición, la recomendación operativa es trabajar dos bloques breves y bien cronometrados, con una pausa mínima entre ambos. Esta estructura permite conservar tres pruebas escritas de 100 minutos sin salir del período disponible.
\end{warnbox}

\section{Contenidos por unidad}

\subsection{Unidad II. Meteorología e hidrología}
\unitbox
{Unidad II}
{La exposición debe conectar la base meteorológica con los procesos hidrológicos que controlan disponibilidad, pérdida y movimiento del agua. No conviene convertirla en una clase general de meteorología aislada de la hidrología.}
{Programa de la asignatura; \textcite{aparicio1989} como apoyo parcial desde nociones de hidrometeorología y precipitación; \textcite{singh2016} para clima, hidrometeorología y cambio climático.}

\textbf{Contenidos mínimos}
\begin{itemize}
\item Definición de elementos climatológicos relevantes para hidrología.
\item Medición y registro de variables meteorológicas.
\item Radiación solar, temperatura, humedad y vientos como controladores de procesos hidrológicos.
\item Evaporación como vínculo entre atmósfera y balance hidrológico.
\item Impactos hidrológicos del cambio climático, al menos a nivel conceptual.
\end{itemize}

\textbf{Contenidos recomendados}
\begin{itemize}
\item Diferencia entre tiempo atmosférico, clima e hidrometeorología.
\item Instrumentos o redes de observación más usuales y limitaciones de los datos.
\item Ejemplos de estaciones meteorológicas e hidroclimatológicas y de las variables que registran.
\item Ejemplos de registros de precipitación, temperatura, humedad relativa, velocidad del viento, evaporación u otras variables discutidas.
\item Relación entre variabilidad climática, extremos y disponibilidad hídrica.
\item Ejemplo breve de cómo una variable meteorológica afecta precipitación, escurrimiento o evapotranspiración.
\end{itemize}

\subsection{Unidad III. Cuencas hidrográficas}
\unitbox
{Unidad III}
{La exposición debe mostrar por qué la cuenca es la unidad natural de análisis hidrológico y cómo sus rasgos fisiográficos condicionan la respuesta frente a la precipitación y la gestión del recurso.}
{Programa de la asignatura; \textcite{aparicio1989}, capítulo sobre cuenca hidrológica; \textcite{mays2011} como apoyo para lectura de respuesta hidrológica y relaciones lluvia--escurrimiento.}

\textbf{Contenidos mínimos}
\begin{itemize}
\item Definición y clasificación de cuencas.
\item Delineación de cuencas por el método New Hampshire.
\item Características fisiográficas relevantes: área, forma, relieve, red de drenaje y cauce principal.
\item Relación entre características fisiográficas y respuesta hidrológica.
\item Manejo de cuencas hidrográficas y su importancia.
\item Herramientas de gestión de recursos hídricos y planificación del manejo de cuencas.
\item Importancia del bosque en la protección de cuencas y deterioro de cuencas hidrográficas.
\end{itemize}

\textbf{Contenidos recomendados}
\begin{itemize}
\item Diferencia entre cuenca superficial y cuenca subterránea.
\item Orden de corrientes, densidad de drenaje y su interpretación.
\item Ejemplo simple de delimitación o lectura de parámetros morfométricos.
\item Conexión entre uso del suelo, erosión, caudales pico y calidad del agua.
\end{itemize}

\subsection{Unidad IV. Precipitación}
\unitbox
{Unidad IV}
{La prioridad es pasar de la formación de la lluvia a la forma en que la precipitación se mide, se representa y se transforma en información útil para análisis hidrológico y diseño.}
{Programa de la asignatura; \textcite{aparicio1989}, capítulo de precipitación; \textcite{chow1994} para apoyo adicional en tratamiento de datos y análisis areal; \textcite{montealegre1990} para consistencia y control de registros.}

\textbf{Contenidos mínimos}
\begin{itemize}
\item Formación de la lluvia.
\item Medición y representación de la precipitación.
\item Estimación de precipitación media en una zona.
\item Deducción o ajuste de datos faltantes en registros.
\item Curvas altura--área--duración.
\item Análisis e interpretación de datos de precipitación.
\end{itemize}

\textbf{Contenidos recomendados}
\begin{itemize}
\item Tipos de precipitación y su relevancia hidrológica.
\item Pluviómetro y pluviógrafo; diferencias entre dato acumulado e intensidad.
\item Métodos areales como promedio aritmético, polígonos de Thiessen o isoyetas.
\item Ejemplo breve de completación o ajuste de un registro.
\item Método de la curva de doble masa para revisión de consistencia y ajuste de registros.
\end{itemize}

\subsection{Unidad V. Escurrimiento}
\unitbox
{Unidad V}
{La exposición debe explicar el origen del escurrimiento, su medición y sus representaciones básicas, con énfasis en cómo se interpreta la respuesta de una cuenca.}
{Programa de la asignatura; \textcite{aparicio1989}, capítulo de escurrimiento; \textcite{chow1994} y \textcite{mays2011} para aforo, curva de gastos, urbanización e hidrogramas unitarios.}

\textbf{Contenidos mínimos}
\begin{itemize}
\item Proceso de escurrimiento y clasificación.
\item Diferencia entre escurrimiento superficial, subsuperficial y subterráneo.
\item Aforo de corrientes superficiales.
\item Representación del escurrimiento.
\item Hidrograma.
\item Hidrogramas unitarios.
\item Curva de gastos.
\end{itemize}

\textbf{Contenidos recomendados}
\begin{itemize}
\item Escurrimiento directo y escurrimiento base.
\item Variables típicas de un hidrograma: ascenso, pico, recesión, volumen.
\item Efectos de la urbanización sobre la forma del hidrograma y sobre el gasto pico.
\item Relación entre aforo, sección, velocidad y gasto.
\item Ejemplo simple de lectura o construcción de curva de gastos.
\end{itemize}

\subsection{Unidad VI. Predicción de eventos de diseño}
\unitbox
{Unidad VI}
{La exposición debe centrarse en cómo los datos hidrológicos se depuran, se interpretan probabilísticamente y se traducen en eventos de diseño para obras civiles.}
{Programa de la asignatura; \textcite{aparicio1989} para frecuencia hidrológica y relaciones lluvia--escurrimiento; \textcite{chow1994} y \textcite{mays2011} para análisis de frecuencia, IDF y distribuciones comunes; \textcite{singh2016} para extremos hidrometeorológicos y contexto de cambio climático.}

\textbf{Contenidos mínimos}
\begin{itemize}
\item Obtención de una muestra de datos y criterios básicos de selección.
\item Extensión de registros.
\item Pruebas de homogeneidad e independencia.
\item Período de retorno.
\item Estimación de gastos máximos anuales.
\item Curvas intensidad--duración--período de retorno.
\item Estimación de gastos mínimos anuales.
\item Mención de modelos probabilísticos de uso común, en particular EV1 (Gumbel) y Log-Pearson Tipo III.
\end{itemize}

\textbf{Contenidos recomendados}
\begin{itemize}
\item Diferencia entre frecuencia empírica y modelo probabilístico.
\item Interpretación correcta del período de retorno y errores comunes.
\item Papel de la longitud y calidad del registro en la confiabilidad del resultado.
\item Ejemplo de lectura de una curva IDF o de comparación entre eventos máximos y mínimos.
\end{itemize}

\subsection{Unidad VII. Evaporación y transpiración}
\unitbox
{Unidad VII}
{La exposición debe mostrar la evaporación y la transpiración como pérdidas hidrológicas centrales del balance de agua, no como fenómenos aislados.}
{Programa de la asignatura; \textcite{aparicio1989}, capítulo de evaporación y transpiración; \textcite{allen1998} para métodos de estimación y uso de Penman--Monteith; \textcite{mays2011} como apoyo.}

\textbf{Contenidos mínimos}
\begin{itemize}
\item Proceso de evaporación y clasificación.
\item Factores que afectan la evaporación.
\item Medición y cálculo de la evaporación.
\item Evapotranspiración.
\item Uso consuntivo.
\item Tipos de modelo para estimación: balance energético, aerodinámico y combinado.
\end{itemize}

\textbf{Contenidos recomendados}
\begin{itemize}
\item Diferencia entre evaporación, transpiración y evapotranspiración potencial o real.
\item Instrumentos o métodos de estimación más usuales.
\item Métodos frecuentes como Penman, Penman--Monteith, tanque de evaporación y relaciones empíricas cuando proceda.
\item Influencia de radiación, viento, temperatura y humedad.
\item Ejemplo breve de aplicación, por ejemplo estimación de evapotranspiración de referencia con FAO-56 o análisis de pérdidas en embalses.
\end{itemize}

\subsection{Unidad VIII. Infiltración}
\unitbox
{Unidad VIII}
{La exposición debe explicar la infiltración como proceso de ingreso de agua al suelo y como control directo de escurrimiento, recarga y pérdidas.}
{Programa de la asignatura; \textcite{aparicio1989}, capítulo de infiltración; \textcite{chow1994} y \textcite{mays2011} para métodos de infiltración y abstracciones.}

\textbf{Contenidos mínimos}
\begin{itemize}
\item Proceso de infiltración.
\item Factores que afectan la infiltración.
\item Métodos para calcular la infiltración.
\item Medición de la infiltración.
\item Referencia explícita a enfoques usuales como Horton, Green--Ampt y Número de Curva del NRCS, aclarando su ámbito de uso.
\end{itemize}

\textbf{Contenidos recomendados}
\begin{itemize}
\item Capacidad de infiltración y su variación temporal.
\item Relación entre infiltración, humedad antecedente, uso del suelo y escurrimiento.
\item Métodos o criterios prácticos como índices de infiltración, Horton, Green--Ampt y Número de Curva del NRCS si se usan con cuidado.
\item Ejemplo simple de interpretación de una curva de infiltración o de un ensayo de campo.
\end{itemize}

\subsection{Unidad IX. Agua subterránea}
\unitbox
{Unidad IX}
{La exposición debe presentar los conceptos fundamentales del agua subterránea con énfasis en recarga, propiedades del acuífero, flujo y respuesta de pozos. En esta unidad, Aparicio no es la fuente principal; debe complementarse de manera más fuerte con otras referencias del curso.}
{Programa de la asignatura; \textcite{freeze1979} como referencia principal para flujo subterráneo, propiedades del acuífero y respuesta de pozos; \textcite{mays2011} como apoyo para integración con ingeniería de recursos hídricos.}

\textbf{Contenidos mínimos}
\begin{itemize}
\item Conceptos de agua subterránea y recarga.
\item Tipos de acuíferos.
\item Porosidad, rendimiento específico y retención específica.
\item Coeficiente de almacenaje, permeabilidad y transmisividad.
\item Ley de Darcy.
\item Pozos y cono de abatimiento.
\end{itemize}

\textbf{Contenidos recomendados}
\begin{itemize}
\item Diferencia entre acuífero libre, confinado y semiconfinado.
\item Superficie freática y gradiente hidráulico.
\item Conexión entre infiltración, percolación y recarga.
\item Ejemplo conceptual de flujo a un pozo o de interpretación del descenso piezométrico.
\end{itemize}

\section{Recomendaciones de cierre para todas las exposiciones}
\begin{itemize}
\item Cerrar con tres o cuatro ideas técnicas que la clase deba retener.
\item Indicar variables que un ingeniero realmente mediría o revisaría en práctica.
\item Explicitar errores frecuentes de interpretación del tema.
\item Conectar la unidad expuesta con otras unidades del curso y con diseño hidrológico.
\end{itemize}

\appendix
\section{Anexo. Formulario breve de audiencia y revisión de pares}
\begin{infobox}
Este formulario es individual. Sirve para controlar asistencia, atención y comprensión mínima de la sesión. Puede ser compartido con el equipo expositor sin identificar al estudiante.
\end{infobox}

\begin{tabularx}{\linewidth}{p{0.28\linewidth} Y}
\toprule
\textbf{Campo} & \textbf{Completar} \\
\midrule
Nombre del estudiante & \\
Matrícula & \\
Fecha & \\
Unidad expuesta & \\
Equipo expositor & \\
\bottomrule
\end{tabularx}

\subsection*{1. Verificación rápida}
\begin{tabularx}{\linewidth}{Y p{0.20\linewidth}}
\toprule
\textbf{Ítem} & \textbf{Respuesta} \\
\midrule
Asistí a la sesión completa de exposición. & Sí / No \\
Tomé nota de al menos una idea o dato técnico relevante. & Sí / No \\
Formulé una pregunta, observación o comentario técnico, o puedo dejarlo por escrito abajo. & Sí / No \\
\bottomrule
\end{tabularx}

\subsection*{2. Idea técnica principal de la sesión}
\vspace{2.4cm}

\subsection*{3. Observación o pregunta técnica para el equipo expositor}
\vspace{2.8cm}

\subsection*{4. Pregunta técnica breve}
\textit{Redactada por el profesor o tomada del núcleo de la exposición. Debe poder responderse en 3 a 5 líneas.}

\vspace{3.0cm}

\subsection*{5. Valoración breve de la exposición}
\begin{tabularx}{\linewidth}{Y p{0.18\linewidth}}
\toprule
\textbf{Criterio} & \textbf{Nivel} \\
\midrule
La exposición fue clara y ordenada. & \makecell[l]{Alta /\\ Media /\\ Baja} \\
El equipo mostró dominio técnico razonable. & \makecell[l]{Alta /\\ Media /\\ Baja} \\
El material visual ayudó a comprender el tema. & \makecell[l]{Alta /\\ Media /\\ Baja} \\
\bottomrule
\end{tabularx}

\begin{warnbox}
La entrega de este formulario sin contenido técnico verificable puede ser registrada como participación insuficiente.
\end{warnbox}

\printbibliography[title={Referencias}]

\end{document}
